\chapter{Addition, negation, and surreal numbers}

\section{Addition and Negation on games}

\begin{definition}[Recursive sum]
  \label{def:game-add}
  \lean{Game.add,add_left,add_right,mem_add_left_iff,mem_add_right_iff}
  \leanok
  \uses{def:game,def:birthday,lem:birthday-options}
  For games \(x=\{X^L\mid X^R\}\) and \(y=\{Y^L\mid Y^R\}\), define
  \[
    x+y=
    \bigl\{\,x^L+y,\ x+y^L\ \bigm|\ x^R+y,\ x+y^R\,\bigr\},
  \]
  with one option for every member of the indicated finite lists.
\end{definition}

\begin{theorem}[Basic addition laws]
  \label{thm:game-add-laws}
  \lean{Game.add_zero',Game.zero_add',Game.add_zero,Game.zero_add,Game.add_comm,Game.add_assoc}
  \leanok
  For all games \(a,b,c\),
  \[
    a+0=a,\qquad 0+a=a,\qquad a+b\GameEq b+a,
    \qquad (a+b)+c=a+(b+c).
  \]
  The zero and associativity identities hold as literal equalities of game trees; commutativity
  is asserted up to game equivalence.
\end{theorem}
\begin{proof}
  \leanok
  \uses{def:game,def:game-equivalence,lem:birthday-options,def:game-inductions,thm:game-order-calculus,thm:option-extensionality,def:game-add}
  For \(a+0=a\), use the one-game induction scheme of
  Definition~\ref{def:game-inductions}.  Write \(a=\{A^L\mid A^R\}\).  By
  Definition~\ref{def:game-add}, the left and right option lists of \(a+0\)
  are respectively
  \[
    [\,u+0:u\in A^L\,],
    \qquad
    [\,v+0:v\in A^R\,].
  \]
  Every option has smaller birthday by Lemma~\ref{lem:birthday-options}, so
  the induction hypothesis identifies these entries literally with \(u\) and
  \(v\).  Thus the two option lists agree with those of \(a\), and hence
  \(a+0=a\) as game trees.  The proof of \(0+a=a\) is the same, with the
  summands interchanged.

  For commutativity, use the two-game induction scheme of
  Definition~\ref{def:game-inductions}.  By Definition~\ref{def:game-add}, a
  left option of \(a+b\) is either \(a^L+b\) or \(a+b^L\).  In the first
  case Lemma~\ref{lem:birthday-options} makes \((a^L,b)\) a smaller pair, and
  the induction hypothesis gives
  \[
    a^L+b\GameEq b+a^L;
  \]
  the right-hand side is a left option of \(b+a\).  In the second case the
  smaller pair \((a,b^L)\) gives
  \[
    a+b^L\GameEq b^L+a.
  \]
  The two right-option cases are identical, and interchanging \(a,b\) gives
  the converse matching.  Theorem~\ref{thm:option-extensionality} now yields
  \(a+b\GameEq b+a\).

  Associativity is stronger: it is literal equality of game trees.  Apply the
  one-game induction scheme of Definition~\ref{def:game-inductions}
  successively to \(a\), \(b\), and \(c\), generalizing the later variables at
  each stage.  The resulting induction hypotheses are
  \[
    \begin{array}{ll}
      ((a^L+b)+c)=a^L+(b+c) & (a^L\in A^L),\\
      ((a+b^L)+c)=a+(b^L+c) & (b^L\in B^L),\\
      ((a+b)+c^L)=a+(b+c^L) & (c^L\in C^L),
    \end{array}
  \]
  and the analogous three identities for right options.  Each recursive call
  is legitimate by Lemma~\ref{lem:birthday-options}; generalizing the later
  coordinates permits the other two games to remain unchanged.

  Expanding Definition~\ref{def:game-add} twice, the left option lists are
  \[
    \begin{aligned}
      ((a+b)+c)^L
       ={}&[\,(a^L+b)+c:a^L\in A^L\,]\\
          &{}\cup[\,(a+b^L)+c:b^L\in B^L\,]\\
          &{}\cup[\,(a+b)+c^L:c^L\in C^L\,],\\[2mm]
      (a+(b+c))^L
       ={}&[\,a^L+(b+c):a^L\in A^L\,]\\
          &{}\cup[\,a+(b^L+c):b^L\in B^L\,]\\
          &{}\cup[\,a+(b+c^L):c^L\in C^L\,].
    \end{aligned}
  \]
  We use \(\cup\) for the union of these finite option collections; formally, the underlying
  lists retain order and multiplicity.  On the left the three sublists are parenthesized as
  \((A^L\cup B^L)\cup C^L\), whereas on the right they are parenthesized as
  \(A^L\cup(B^L\cup C^L)\).  Associativity of finite-list concatenation first identifies these
  two parenthesizations.
  The lists then contain the same three families in the same order; only their
  entries differ.  The first, second, and third induction
  hypotheses above give pointwise equality on \(A^L,B^L,C^L\), respectively.
  Consequently the three corresponding sublists are equal, and therefore the
  complete left lists are equal.  Replacing \(A^L,B^L,C^L\) by
  \(A^R,B^R,C^R\) proves equality of the right lists in the same way.  The two
  games have identical left and right option lists, so
  \(((a+b)+c)=a+(b+c)\) literally.
\end{proof}
\begin{theorem}[Translation and monotonicity]
  \label{thm:game-add-order}
  \lean{Game.add_le_add_right,Game.add_right_cancel,Game.add_le_add,Game.add_equal,Game.add_lt_le,Game.add_le_lt,Game.add_lt_left_left,Game.add_lt_right_right}
  \leanok
  Addition preserves and reflects weak order in either variable.  In particular,
  \[
    a\GameLe b\quad\Longleftrightarrow\quad a+c\GameLe b+c,
  \]
  and
  \[
    a\GameLe c,\ b\GameLe d\Longrightarrow a+b\GameLe c+d.
  \]
  If either input comparison is strict and the other weak, the resulting comparison is strict.
  Addition also respects \(\GameEq\), i.e.
  \[
    a\GameEq b\quad\Longleftrightarrow\quad a+c\GameEq b+c.
  \]
\end{theorem}
\begin{proof}
  \leanok
  \uses{lem:birthday-options,def:game-order,def:game-equivalence,def:game-inductions,thm:game-order-calculus,def:game-add,thm:game-add-laws}
  We first prove simultaneously, in both directions, the equivalence
  \[
    a+c\GameLe b+c\quad\Longleftrightarrow\quad a\GameLe b.
  \]
  Use the three-game induction scheme of Definition~\ref{def:game-inductions}, with measure
  \(\bd(a)+\bd(b)+\bd(c)\), and use the comparison of
  Definition~\ref{def:game-order} throughout.

  Suppose first that \(a+c\GameLe b+c\).  If \(a^L\) is a left option of
  \(a\) and \(b\GameLe a^L\), the induction hypothesis for
  \((b,a^L,c)\) gives \(b+c\GameLe a^L+c\).  Transitivity from
  Theorem~\ref{thm:game-order-calculus}, together with the assumed comparison,
  gives \(a+c\GameLe a^L+c\).  This is impossible because
  Definition~\ref{def:game-add} makes \(a^L+c\) a left option of \(a+c\),
  while the left-option clause of the reflexive comparison
  \(a+c\GameLe a+c\) forbids such an inequality.  If \(b^R\) is a right
  option of \(b\) and \(b^R\GameLe a\), the induction hypothesis gives
  \(b^R+c\GameLe a+c\); composing with the assumed comparison contradicts the
  right-option clause of the reflexive comparison \(b+c\GameLe b+c\).  Thus
  \(a\GameLe b\).

  Conversely, assume \(a\GameLe b\).  By Definition~\ref{def:game-add}, the
  left options of \(a+c\) are \(a^L+c\) and \(a+c^L\), while the right
  options of \(b+c\) are \(b^R+c\) and \(b+c^R\).  If
  \(b+c\GameLe a^L+c\), the induction hypothesis for \((b,a^L,c)\) gives
  \(b\GameLe a^L\), contrary to the left-option clause of \(a\GameLe b\).
  If \(b+c\GameLe a+c^L\), the induction hypothesis for \((a,b,c^L)\)
  gives \(a+c^L\GameLe b+c^L\).  Transitivity would then give
  \(b+c\GameLe b+c^L\), contradicting the left-option clause of the reflexive
  comparison \(b+c\GameLe b+c\).  The right-option cases are dual.  A
  comparison \(b^R+c\GameLe a+c\) reflects to \(b^R\GameLe a\), which is
  forbidden by \(a\GameLe b\).  Finally, if
  \(b+c^R\GameLe a+c\), the induction hypothesis for \((a,b,c^R)\) gives
  \(a+c^R\GameLe b+c^R\); transitivity would imply
  \(a+c^R\GameLe a+c\), contrary to the right-option clause of the reflexive
  comparison \(a+c\GameLe a+c\).  This proves the reverse implication in the
  displayed equivalence.
  Every recursive call decreases because the coordinate replaced in it is an
  option, hence has smaller birthday by Lemma~\ref{lem:birthday-options}.

  Commutativity from Theorem~\ref{thm:game-add-laws} turns the displayed
  equivalence into the
  corresponding equivalence for addition on the left.  If
  \(a\GameLe c\) and \(b\GameLe d\), change first one summand and then the
  other, and compose the two inequalities using
  Theorem~\ref{thm:game-order-calculus}; this proves two-variable
  monotonicity.  Suppose, for example, that \(a\GameLt c\) and
  \(b\GameLe d\).  Weak monotonicity gives \(a+b\GameLe c+d\).  If the reverse
  comparison held, then
  \[
    c+b\GameLe c+d\GameLe a+b.
  \]
  Cancelling the common summand \(b\) by the displayed translation equivalence
  would give \(c\GameLe a\), contradicting \(a\GameLt c\).  Thus
  \(a+b\GameLt c+d\).  Strictness in the second variable is proved
  symmetrically.  Applying the weak result in both directions proves
  compatibility with \(\GameEq\).
\end{proof}

\begin{definition}[Negation]
  \label{def:game-neg}
  \lean{Game.neg,neg_left_def,neg_right_def,Game.neg_zero}
  \leanok
  \uses{def:game,def:birthday,lem:birthday-options}
  Negation interchanges the two option lists and recursively negates every option:
  \[
    -\{X^L\mid X^R\}=\{-X^R\mid -X^L\}.
  \]
  In particular, \(-0=0\) as a literal equality of game trees.
\end{definition}

\begin{theorem}[Negation reverses order]
  \label{thm:game-neg-order}
  \lean{Game.neg_le_neg,Game.neg_congr,Game.neg_lt_neg}
  \leanok
  For all games \(a,b\),
  \[
    a\GameLe b\Longleftrightarrow -b\GameLe-a,
    \qquad
    a\GameLt b\Longleftrightarrow -b\GameLt-a,
    \qquad
    a\GameEq b\Longleftrightarrow -b\GameEq-a.
  \]
\end{theorem}
\begin{proof}
  \leanok
  \uses{lem:birthday-options,def:game-order,def:game-equivalence,def:game-inductions,def:game-neg}
  Use the two-game induction scheme of Definition~\ref{def:game-inductions} to
  prove
  \[
    a\GameLe b\quad\Longleftrightarrow\quad -b\GameLe-a.
  \]
  Assume first \(a\GameLe b\).  By Definition~\ref{def:game-neg}, a left
  option of \(-b\) has the form \(-b^R\), and a right option of \(-a\) has
  the form \(-a^L\).  If \(-a\GameLe-b^R\), the induction hypothesis for
  \((b^R,a)\) gives \(b^R\GameLe a\), contradicting the right-option clause
  of \(a\GameLe b\) in Definition~\ref{def:game-order}.  Likewise, if
  \(-a^L\GameLe-b\), the induction hypothesis for \((b,a^L)\) gives
  \(b\GameLe a^L\), contradicting the left-option clause.  These are exactly
  the two clauses required for \(-b\GameLe-a\).

  Conversely, assume \(-b\GameLe-a\).  To verify \(a\GameLe b\), let
  \(a^L\) be a left option of \(a\).  A comparison \(b\GameLe a^L\) would,
  by the induction hypothesis for \((b,a^L)\), imply
  \(-a^L\GameLe-b\), contradicting the right-option clause of
  \(-b\GameLe-a\).  The argument for a right option \(b^R\) is the same:
  \(b^R\GameLe a\) would imply \(-a\GameLe-b^R\), contradicting the
  left-option clause.  Each induction step is valid by
  Lemma~\ref{lem:birthday-options}.

  Applying the displayed equivalence to both weak inequalities in
  Definition~\ref{def:game-equivalence} proves the stated biconditional for equivalence.
  Applying it to the forward weak
  inequality and to the excluded reverse inequality proves the strict-order
  equivalence.  Both conclusions use
  Definitions~\ref{def:game-order} and~\ref{def:game-equivalence}.
\end{proof}

\begin{theorem}[Additive inverses]
  \label{thm:game-add-inverse}
  \lean{Game.add_neg,Game.neg_add}
  \leanok
  Every game satisfies
  \[
    x+(-x)\GameEq0,
    \qquad
    (-x)+x\GameEq0.
  \]
\end{theorem}
\begin{proof}
  \leanok
  \uses{def:game,def:game-order,def:game-equivalence,def:game-inductions,lem:birthday-options,thm:game-order-calculus,def:game-add,thm:game-add-laws,def:game-neg}
  We argue by induction on \(\bd(x)\), using the one-game scheme of
  Definition~\ref{def:game-inductions}.  Put \(s=x+(-x)\).  We prove the two
  inequalities in the definition of \(s\GameEq0\) from
  Definition~\ref{def:game-equivalence}.  Since \(0=\{\,\mid\,\}\) by
  Definition~\ref{def:game}, the clauses involving options of \(0\) are
  vacuous.

  By Definitions~\ref{def:game-add} and~\ref{def:game-neg}, a left option of
  \(s\) is either
  \[
    \ell=x^L+(-x)
    \qquad\text{or}\qquad
    \ell=x+(-x^R).
  \]
  In the first case \(x^L+(-x^L)\) is a right option of \(\ell\); in the
  second case \(x^R+(-x^R)\) is a right option of \(\ell\).  By
  Lemma~\ref{lem:birthday-options} and the induction hypothesis, the indicated
  self-cancelling game is equivalent to \(0\), and in particular is weakly
  below \(0\).  Therefore \(0\GameLe\ell\) would violate the right-option
  clause of that comparison.  This verifies the left-option clause of
  \(s\GameLe0\).

  Similarly, a right option of \(s\) is either
  \[
    r=x^R+(-x)
    \qquad\text{or}\qquad
    r=x+(-x^L).
  \]
  The first has \(x^R+(-x^R)\) as a left option, and the second has
  \(x^L+(-x^L)\) as a left option.  The induction hypothesis makes that left
  option weakly above \(0\).  Thus \(r\GameLe0\) would violate its
  left-option clause.  This proves \(0\GameLe s\), and hence
  \(x+(-x)\GameEq0\).  Commutativity from
  Theorem~\ref{thm:game-add-laws} gives \((-x)+x\GameEq0\).
\end{proof}

\section{Surreal Number}

\begin{theorem}[Surreal games are closed under addition and negation]
  \label{thm:surreal-add-closure}
  \lean{Surreal.add_isSurreal,Surreal.neg_isSurreal,Surreal.add,Surreal.neg}
  \leanok
  If \(a,b\) are surreal representatives, then the games \(a+b\) and \(-a\) are surreal.
  Hence addition and negation define operations on surreal representatives.
\end{theorem}
\begin{proof}
  \leanok
  \uses{def:birthday,def:surreal-pair-induction,lem:birthday-options,def:game-order,def:is-surreal,thm:game-order-calculus,thm:between-options,def:surreal-representative,def:game-add,thm:game-add-order,def:game-neg,thm:game-neg-order}
  For addition, use the well-founded induction on pairs of surreal
  representatives from Definition~\ref{def:surreal-pair-induction}.  By
  Definition~\ref{def:game-add}, every option of \(a+b\) replaces exactly one
  of \(a,b\) by one of its options.  Lemma~\ref{lem:birthday-options} makes
  the resulting pair smaller, so the induction hypothesis proves that every
  option of \(a+b\) is surreal, as required by
  Definition~\ref{def:is-surreal}.

  It remains to separate a left option from a right option.  There are four
  cases:
  \[
    \begin{array}{c|c}
      \text{left option}&\text{right option}\\ \hline
      a^L+b&a^R+b\\
      a+b^L&a+b^R\\
      a^L+b&a+b^R\\
      a+b^L&a^R+b.
    \end{array}
  \]
  In the first case Theorem~\ref{thm:between-options} and strict transitivity
  from Theorem~\ref{thm:game-order-calculus} give
  \(a^L\GameLt a^R\), and strict translation invariance from
  Theorem~\ref{thm:game-add-order} gives
  \(a^L+b\GameLt a^R+b\).  The second case is identical with \(a,b\)
  interchanged.  In the third case,
  \(a^L\GameLt a\) and \(b\GameLe b^R\); in the fourth,
  \(a\GameLt a^R\) and \(b^L\GameLe b\).  The two mixed monotonicity
  statements of Theorem~\ref{thm:game-add-order} again put the chosen left
  option strictly below the chosen right option.  By
  Definition~\ref{def:game-order}, this strict inequality excludes the reverse
  weak comparison, which is exactly the separation condition in
  Definition~\ref{def:is-surreal}.

  For negation, induct on the birthday of the underlying game of the surreal
  representative \(a\).  Definition~\ref{def:game-neg} says that
  a left option of \(-a\) is \(-a^R\), while a right option is \(-a^L\).
  Theorem~\ref{thm:between-options} and strict transitivity from
  Theorem~\ref{thm:game-order-calculus} give
  \(a^L\GameLt a\GameLt a^R\), hence \(a^L\GameLt a^R\).  If
  \(-a^L\GameLe-a^R\), order reversal in
  Theorem~\ref{thm:game-neg-order} would give \(a^R\GameLe a^L\), contradicting
  that strict comparison.  Thus the left and right options of \(-a\) satisfy
  the separation condition of Definition~\ref{def:is-surreal}.
  Each negated option is surreal by the induction hypothesis and
  Lemma~\ref{lem:birthday-options}.  Thus \(-a\) satisfies
  Definition~\ref{def:is-surreal}.  Finally, Definition~\ref{def:surreal-representative}
  turns each game, together with the surrealness just proved, into the
  corresponding operation on representatives.
\end{proof}

\begin{definition}[Auxiliary quotient of all games]
  \label{def:game-quotient}
  \lean{Game.instSetoidGame,Game.GameQ,Game.q_eq}
  \leanok
  \uses{def:game-equivalence,thm:game-order-calculus}
  By Theorem~\ref{thm:game-order-calculus}, game equivalence is an equivalence relation.
  Define \(\Games/{\GameEq}\) to be the quotient of \(\Games\) by game equivalence, and
  write \([a]\) for the equivalence class of a game \(a\).  Thus \([a]=[b]\) if and only
  if \(a\GameEq b\).
\end{definition}

\begin{theorem}[Partially ordered additive commutative group]
  \label{thm:game-quotient-ordered-add-group}
  \lean{Game.q_zero,Game.q_add,Game.q_neg,Game.q_le,Game.q_lt,Game.instAddCommGroupGameQ,Game.instPartialOrderGameQ,Game.instIsOrderedAddMonoidGameQ}
  \leanok
  For finite games \(a,b\), define
  \[
    0=[0],\qquad [a]+[b]=[a+b],\qquad -[a]=[-a],
  \]
  and
  \[
    [a]\leq[b]\quad\Longleftrightarrow\quad a\GameLe b,
    \qquad
    [a]<[b]\quad\Longleftrightarrow\quad a\GameLt b.
  \]
  These formulas are independent of the chosen representatives.  With these operations and
  this order, \(\Games/{\GameEq}\) is a partially ordered additive commutative group, and
  addition is order preserving in each variable.
\end{theorem}
\begin{proof}
  \leanok
  \uses{def:game-quotient,def:game-order,def:game-equivalence,def:game-add,def:game-neg,thm:game-order-calculus,thm:game-add-laws,thm:game-add-order,thm:game-neg-order,thm:game-add-inverse}
  The addition and negation formulas are independent of the representatives by the
  congruence assertions in Theorems~\ref{thm:game-add-order}
  and~\ref{thm:game-neg-order}.

  The order is also independent of representatives.  Indeed, suppose
  \(a\GameEq a'\), \(b\GameEq b'\), and \(a\GameLe b\).  The two appropriate
  components of the equivalences give \(a'\GameLe a\) and
  \(b\GameLe b'\); transitivity from
  Theorem~\ref{thm:game-order-calculus} yields \(a'\GameLe b'\).  The converse
  follows symmetrically.  Reflexivity and transitivity therefore pass to
  classes.  If \([a]\leq[b]\) and \([b]\leq[a]\), then
  \(a\GameEq b\) by Definition~\ref{def:game-equivalence}, so \([a]=[b]\) by the
  definition of an equivalence class.  Thus the quotient order is a partial
  order.  The asserted characterization of strict order now follows from the
  definitions of \(\GameLt\) and strict order in a partial order.

  The zero, associativity, and commutativity laws descend from
  Theorem~\ref{thm:game-add-laws}, and the inverse law descends from
  Theorem~\ref{thm:game-add-inverse}; an equivalence of representatives is
  precisely equality of their classes.  Finally, translation monotonicity is
  the class-level form of Theorem~\ref{thm:game-add-order}.  Commutativity from
  Theorem~\ref{thm:game-add-laws} permits the common summand to be placed on
  either side.
\end{proof}

\begin{definition}[Short surreal numbers]
  \label{def:surreal-number}
  \lean{Surreal.SurrealNumber,Surreal.SurrealNumber.add,Surreal.SurrealNumber.neg}
  \leanok
  \uses{def:surreal-representative,thm:surreal-add-closure,thm:game-order-calculus,thm:game-add-order,thm:game-neg-order}
  The type \(\NoShort\) is the quotient of surreal representatives by \(a\GameEq b\).
  Addition, zero, one, and negation descend to this quotient.
\end{definition}

\begin{theorem}[Linear ordered additive commutative group]
  \label{thm:ordered-add-group}
  \lean{Surreal.instLinearOrderSurrealNumber,Surreal.instAddCommGroupSurrealNumber,Surreal.instIsOrderedAddMonoidSurrealNumber}
  \leanok
  The quotient \(\NoShort\) is a linearly ordered additive commutative group, and addition is
  compatible with its order.
\end{theorem}
\begin{proof}
  \leanok
  \uses{def:surreal-representative,thm:surreal-totality,thm:surreal-add-closure,def:game-quotient,thm:game-quotient-ordered-add-group,def:surreal-number}
  There is a natural injective map
  \[
    \iota:\NoShort\longrightarrow\Games/{\GameEq},
    \qquad \iota([s])=[s_{\mathrm{game}}].
  \]
  This is well defined and injective because both quotients identify representatives exactly
  when their underlying games are game equivalent.  Pull back the quotient order along
  \(\iota\), so \([s]\leq[t]\) precisely when
  \(s_{\mathrm{game}}\GameLe t_{\mathrm{game}}\).  Definition~\ref{def:surreal-number} and
  Theorem~\ref{thm:surreal-add-closure} show that the image contains zero and is closed under
  addition and negation.  Hence Theorem~\ref{thm:game-quotient-ordered-add-group} supplies the
  additive commutative group, partial order, and order compatibility on \(\NoShort\).
  Theorem~\ref{thm:surreal-totality} upgrades this partial order to a linear order.
\end{proof}
